\uuid{7ErJ}
\exo7id{719}
\titre{exo7 719}
\auteur{bodin}
\organisation{exo7}
\datecreate{1998-09-01}
\isIndication{false}
\isCorrection{false}
\chapitre{Dérivabilité des fonctions réelles}
\sousChapitre{Théorème de Rolle et accroissements finis}
\cpge{mpsi-mp2i.nombres-complexes.a, pcsi.nombres-complexes.a, ptsi.nombres-complexes.a, mpsi-mp2i.calcul-algebrique-trigonometrie.c, pcsi.calcul-algebrique-trigonometrie.c, ptsi.calcul-algebrique-trigonometrie.c, bcpst1.nombres-reels, bcpst1.methodes-calcul}
\module{Analyse}
\niveau{L1}
\difficulte{}

\contenu{
\question{
\'Etant donn\'e $y$ un r\'eel positif
et $n$ un entier naturel pair, montrer que
$(x+y)^n=x^n+y^n$ si et seulement si $x=0$.
Cas $n$ impair ?
}
}
